% GENERATED FILE. Edit canonical lessons, then run npm run generate:books.
% canonical-source-hash: fnv1a64:278e634f42e2795f
% canonical-lessons: JA-C301-takara, JA-C301-horu, JA-C301-umeru, JA-C301-binbou, JA-C301-yokubari, JA-R301-review-297-301

\chapter{Buried Treasure}
\label{ch:ja-buried-treasure}
\hlchaptermodality{301}

\begin{chapteropening}
\textbf{By the end of this chapter:} \emph{I can say who buried the treasure and who dug it up.}

Complete the last lesson of chapter 301: I can say who buried the treasure and who dug it up.
\end{chapteropening}

\section[\texorpdfstring{\ja{たから} (takara)}{たから (takara)}]{\ja{たから} (takara) — a treasure}

\label{lesson:JA-C301-takara}



\begin{quote}
Before the new one: say the Japanese for a curse, then the Japanese for sly.
\end{quote}

\subsection*{You'll want to know: \ja{たから}}
\textbf{\ja{たから}} — \emph{takara} — \textquotedblleft{}a treasure\textquotedblright{}.

\begin{quote}
\textbf{\ja{むかし} \ja{ここに} \ja{たからが} \ja{ありました。}} — \emph{mukashi koko ni takara ga arimashita} — \textquotedblleft{}Long ago there was treasure here.\textquotedblright{}
\end{quote}

\begin{cousinweb}
\textbf{\ja{たから}} (\emph{takara}) is a native Japanese word (wago), inherited rather than borrowed from Chinese; its kanji was attached later as a kun'yomi (JMdict).
\end{cousinweb}

\begin{grammarlens}[title={What you've built}]
One more word for telling what happened.
\end{grammarlens}

\subsection*{Your turn}
\begin{itemize}
\raggedright
  \item \emph{Say it:} \emph{takara}
  \item \emph{Say it:} \emph{takara}, once more
  \item \emph{Say it:} \emph{mukashi koko ni takara ga arimashita} — Long ago there was treasure here.
  \item \emph{Recall:} say the Japanese for a curse, then the Japanese for sly, then say \emph{takara} again
\end{itemize}

\begin{tcolorbox}[breakable,colback=teal!4,colframe=teal!35!black,title={Before you move on}]
What does \ja{たから} mean? (\textquotedblleft{}A treasure\textquotedblright{}.) Say it once more.
\end{tcolorbox}
\section[\texorpdfstring{\ja{ほる} (horu)}{ほる (horu)}]{\ja{ほる} (horu) — to dig}

\label{lesson:JA-C301-horu}



\begin{quote}
Before the new one: say the Japanese for sly, then the Japanese for a treasure.
\end{quote}

\subsection*{You'll want to know: \ja{ほる}}
\textbf{\ja{ほる}} — \emph{horu} — \textquotedblleft{}to dig\textquotedblright{}.

\begin{quote}
\textbf{\ja{いぬが} \ja{にわを} \ja{ほっています。}} — \emph{inu ga niwa o hotte imasu} — \textquotedblleft{}The dog is digging in the garden.\textquotedblright{}
\end{quote}

Polite \textbf{\ja{ほります}}, \emph{horimasu}; a godan verb. Past (polite): \textbf{\ja{おじいさんが} \ja{にわを} \ja{ほりました。}}, \emph{ojiisan ga niwa o horimashita}.

\begin{cousinweb}
\textbf{\ja{ほる}} (\emph{horu}) is a native Japanese word (wago), inherited rather than borrowed from Chinese; its kanji was attached later as a kun'yomi (JMdict).
\end{cousinweb}

\begin{grammarlens}[title={What you've built}]
One more word for telling what happened.
\end{grammarlens}

\subsection*{Your turn}
\begin{itemize}
\raggedright
  \item \emph{Say it:} \emph{horu}
  \item \emph{Say it:} \emph{horu}, once more
  \item \emph{Say it:} \emph{inu ga niwa o hotte imasu} — The dog is digging in the garden.
  \item \emph{Recall:} say the Japanese for sly, then the Japanese for a treasure, then say \emph{horu} again
\end{itemize}

\begin{tcolorbox}[breakable,colback=teal!4,colframe=teal!35!black,title={Before you move on}]
What does \ja{ほる} mean? (\textquotedblleft{}To dig\textquotedblright{}.) Say it once more.
\end{tcolorbox}
\section[\texorpdfstring{\ja{うめる} (umeru)}{うめる (umeru)}]{\ja{うめる} (umeru) — to bury}

\label{lesson:JA-C301-umeru}



\begin{quote}
Before the new one: say the Japanese for a treasure, then the Japanese for to dig.
\end{quote}

\subsection*{You'll want to know: \ja{うめる}}
\textbf{\ja{うめる}} — \emph{umeru} — \textquotedblleft{}to bury\textquotedblright{}.

\begin{quote}
\textbf{\ja{おとこは} \ja{たからを} \ja{やまに} \ja{うめました。}} — \emph{otoko wa takara o yama ni umemashita} — \textquotedblleft{}The man buried the treasure in the mountains.\textquotedblright{}
\end{quote}

Polite \textbf{\ja{うめます}}, \emph{umemasu}; an ichidan verb. Past (polite): \textbf{\ja{とのさまが} \ja{たからを} \ja{うめました。}}, \emph{tonosama ga takara o umemashita}.

\begin{cousinweb}
\textbf{\ja{うめる}} (\emph{umeru}) is a native Japanese word (wago), inherited rather than borrowed from Chinese; its kanji was attached later as a kun'yomi (JMdict).
\end{cousinweb}

\begin{grammarlens}[title={What you've built}]
One more word for telling what happened.
\end{grammarlens}

\subsection*{Your turn}
\begin{itemize}
\raggedright
  \item \emph{Say it:} \emph{umeru}
  \item \emph{Say it:} \emph{umeru}, once more
  \item \emph{Say it:} \emph{otoko wa takara o yama ni umemashita} — The man buried the treasure in the mountains.
  \item \emph{Recall:} say the Japanese for a treasure, then the Japanese for to dig, then say \emph{umeru} again
\end{itemize}

\begin{tcolorbox}[breakable,colback=teal!4,colframe=teal!35!black,title={Before you move on}]
What does \ja{うめる} mean? (\textquotedblleft{}To bury\textquotedblright{}.) Say it once more.
\end{tcolorbox}
\section[\texorpdfstring{\ja{びんぼう} (binbō)}{びんぼう (binbō)}]{\ja{びんぼう} (binbō) — poor; poverty}

\label{lesson:JA-C301-binbou}



\begin{quote}
Before the new one: say the Japanese for to dig, then the Japanese for to bury.
\end{quote}

\subsection*{You'll want to know: \ja{びんぼう}}
\textbf{\ja{びんぼう}} — \emph{binbō} — \textquotedblleft{}poor; poverty\textquotedblright{}.

\begin{quote}
\textbf{\ja{その} \ja{おじいさんは} \ja{びんぼうでした。}} — \emph{sono ojiisan wa binbō deshita} — \textquotedblleft{}The old man was poor.\textquotedblright{}
\end{quote}

A na-adjective: \ja{な} before a noun.

\begin{cousinweb}
\textbf{\ja{びんぼう}} (\emph{binbō}) is Sino-Japanese (kango): two Chinese characters read with their on'yomi, readings borrowed from Chinese. They are \emph{bin}, \textquotedblleft{}poor\textquotedblright{} (Middle Chinese \emph{bin}) and \emph{bō}, \textquotedblleft{}scarce\textquotedblright{} (\emph{bjop}) (KANJIDIC2; Middle Chinese in Baxter's transcription).
\end{cousinweb}

\begin{grammarlens}[title={What you've built}]
One more word for telling what happened.
\end{grammarlens}

\subsection*{Your turn}
\begin{itemize}
\raggedright
  \item \emph{Say it:} \emph{binbō}
  \item \emph{Say it:} \emph{binbō}, once more
  \item \emph{Say it:} \emph{sono ojiisan wa binbō deshita} — The old man was poor.
  \item \emph{Recall:} say the Japanese for to dig, then the Japanese for to bury, then say \emph{binbō} again
\end{itemize}

\begin{tcolorbox}[breakable,colback=teal!4,colframe=teal!35!black,title={Before you move on}]
What does \ja{びんぼう} mean? (\textquotedblleft{}Poor; poverty\textquotedblright{}.) Say it once more.
\end{tcolorbox}
\section[\texorpdfstring{\ja{よくばり} (yokubari)}{よくばり (yokubari)}]{\ja{よくばり} (yokubari) — greedy; a greedy person}

\label{lesson:JA-C301-yokubari}



\begin{quote}
Before the new one: say the Japanese for to bury, then the Japanese for poor.
\end{quote}

\subsection*{You'll want to know: \ja{よくばり}}
\textbf{\ja{よくばり}} — \emph{yokubari} — \textquotedblleft{}greedy; a greedy person\textquotedblright{}.

\begin{quote}
\textbf{\ja{となりの} \ja{おじいさんは} \ja{よくばりでした。}} — \emph{tonari no ojiisan wa yokubari deshita} — \textquotedblleft{}The old man next door was greedy.\textquotedblright{}
\end{quote}

A na-adjective: \ja{な} before a noun.

\begin{cousinweb}
\textbf{\ja{よくばり}} (\emph{yokubari}) is a mixed word, native and Sino-Japanese: \emph{yoku}, \textquotedblleft{}greed\textquotedblright{}, is a Chinese character read with its on'yomi, a reading borrowed from Chinese (Middle Chinese \emph{yowk}), and \emph{bari}, \textquotedblleft{}stretch\textquotedblright{}, is native Japanese (JMdict, KANJIDIC2; Middle Chinese in Baxter's transcription).
\end{cousinweb}

\begin{grammarlens}[title={What you've built}]
One more word for telling what happened.
\end{grammarlens}

\subsection*{Your turn}
\begin{itemize}
\raggedright
  \item \emph{Say it:} \emph{yokubari}
  \item \emph{Say it:} \emph{yokubari}, once more
  \item \emph{Say it:} \emph{tonari no ojiisan wa yokubari deshita} — The old man next door was greedy.
  \item \emph{Recall:} say the Japanese for to bury, then the Japanese for poor, then say \emph{yokubari} again
\end{itemize}

\begin{tcolorbox}[breakable,colback=teal!4,colframe=teal!35!black,title={Before you move on}]
What does \ja{よくばり} mean? (\textquotedblleft{}Greedy; a greedy person\textquotedblright{}.) Say it once more.
\end{tcolorbox}
\section[(dialogue)]{Looking back: from The Ogre at the Door to Buried Treasure}

\label{lesson:JA-R301-review-297-301}



\begin{quote}
Say the Japanese for an ogre and a mask.
\end{quote}

\subsection*{Your turn}
Say these aloud:

\begin{itemize}
\raggedright
  \item an ogre, a mask, to be afraid, to drive out, rare
  \item long ago, a folk tale, to tell (a story), a village, a hard worker
  \item a king, a princess, a samurai, to protect, a feudal lord
  \item a fox, a raccoon dog, to take the form of, a curse, sly
  \item a treasure, to dig, to bury, poor, greedy
\end{itemize}

\begin{tcolorbox}[breakable,colback=teal!4,colframe=teal!35!black,title={Before you move on}]
An ogre? (\textbf{\ja{おに}}, \emph{oni}.) A mask? (\textbf{\ja{おめん}}, \emph{omen}.) To be afraid? (\textbf{\ja{こわがる}}, \emph{kowagaru}.) To drive out? (\textbf{\ja{おいだす}}, \emph{oidasu}.) Rare? (\textbf{\ja{めずらしい}}, \emph{mezurashii}.) Long ago? (\textbf{\ja{むかし}}, \emph{mukashi}.) A folk tale? (\textbf{\ja{むかしばなし}}, \emph{mukashibanashi}.) To tell (a story)? (\textbf{\ja{かたる}}, \emph{kataru}.) A village? (\textbf{\ja{むら}}, \emph{mura}.) A hard worker? (\textbf{\ja{はたらきもの}}, \emph{hatarakimono}.) A king? (\textbf{\ja{おうさま}}, \emph{ōsama}.) A princess? (\textbf{\ja{おひめさま}}, \emph{ohimesama}.) A samurai? (\textbf{\ja{さむらい}}, \emph{samurai}.) To protect? (\textbf{\ja{まもる}}, \emph{mamoru}.) A feudal lord? (\textbf{\ja{とのさま}}, \emph{tonosama}.) A fox? (\textbf{\ja{きつね}}, \emph{kitsune}.) A raccoon dog? (\textbf{\ja{たぬき}}, \emph{tanuki}.) To take the form of? (\textbf{\ja{ばける}}, \emph{bakeru}.) A curse? (\textbf{\ja{のろい}}, \emph{noroi}.) Sly? (\textbf{\ja{ずるい}}, \emph{zurui}.) A treasure? (\textbf{\ja{たから}}, \emph{takara}.) To dig? (\textbf{\ja{ほる}}, \emph{horu}.) To bury? (\textbf{\ja{うめる}}, \emph{umeru}.) Poor? (\textbf{\ja{びんぼう}}, \emph{binbō}.) Greedy? (\textbf{\ja{よくばり}}, \emph{yokubari}.)
\end{tcolorbox}
